\newglossaryentry{ipc}
{
	name=indice des prix à la consommation,
	description={Mesure statistique permettant d'estimer, entre deux périodes données, la variation moyenne des prix des produits et des services consommés par les ménages, à qualité constante.
	}
}
\newglossaryentry{désutilité marginale du travail}{
	name=désutilité marginale du travail,
	description={Insatisfaction supplémentaire en termes d'utilité, apportée par l'acquisition d'une unité de travail.
	}
}
\newglossaryentry{digeste}{
	name=Digeste,
	description={[\emph{Nom propre}] Recueil méthodique d'extraits des opinions et sentences des juristes romains, réunis sur l'ordre de l'empereur Justinien en 1530 (publié le 16 décembre 1533). Fondamental dans la compréhension historique des idées juridiques, le Digeste forme la deuxième partie du \emph{Corpus iuris civilis}, et se composera de 9000 fragments rédigés par 40 juristes, affirmant certains aspects de la procédure archaïque, du droit des choses, des actions et des obligations, mais également de la possesion, de l'\emph{animus}, et de la distinction entre le droit écrit et non écrit.
	}
}
\newglossaryentry{courbe d'indifférence}{
	name=courbe d'indifférence,
	description={Une courbe d'indifférence est le lieu géométrique (une courbe dans le plan) des paniers entre lesquels l'individu est exactement indifférent. La courbe d'indifférence associée à $x$ est définie comme l'ensemble des paniers y donnant le même niveau de satisfaction que le panier x :
$ \lbrace y/y \backsim x \rbrace $.
	}
}
\newglossaryentry{hard fork}{
	name=hard fork,
	description={En bases de données décentralisées, la notion de \emph{hard fork} consiste en un changement du paradigme de validation d'ensembles de transactions, partagé par tous les noeuds du réseau, pouvant conduire à une division permanente de la chaîne de blocs en cas de concurrence entre différentes versions logicielles.
	}
}
\newglossaryentry{convertibilité}{
	name=convertibilité,
	description={Se définie comme étant la propriété que possède une monnaie nationale d'être librement échangée, la convertibilité s'opère, selon les époques, soit contre de l'or dans le pays d'émission (ou de l'argent dans certains pays), soit contre des monnaies étrangères.
	}
}
\newglossaryentry{taux d'escompte}{
	name=taux d'escompte,
	description={Le taux d'escompte est un taux d'intérêt utilisé sur le marché monétaire, pour les prêts à très court terme (quelques jours).
	}
}
\newglossaryentry{thésaurisation}{
	name=thésaurisation,
	description={Comportement économique résultant de la volonté de conserver la valeur (monnaie, billets, pierres précieuses, or...) en dehors du système économique pour en tirer un profit ou par absence de meilleur emploi.
	}
}
\newglossaryentry{Loi de Gresham}{
	name=Loi de Gresham appliquée au bimétallisme or-argent,
	description={La plus importante rareté relative de l'argent par rapport à l'or, suite à la découverte d'importants gisements en Californie et en Australie dès 1848, conduit à un dépassement de la valeur légale de l'argent par sa valeur physique. Les agents économiques se retrouvent dans une situation ou il devient préférable d'acheter de l'argent avec de l'or au cours légal, afin de le revendre sur les marchés, la banque centrale réglant la différence de valeur par la liquidation de ses réserves d'argent. La reconstitution des réserve entraine un mécanisme plus complexe : la banque centrale se trouve en une situation contraignante, et l'acquisition impérieuse d'argent au prix des marchés conduit les agents à spéculer sur sa réévaluation. Le gain à l'échange couvrant les coûts de refonte, de frappe et de transport, la demande de thésaurisation des monnaies d'argent entrainera à moyen terme une crise des règlements par la diminution relative de circulation des signes monétaires.
	}
}
\newglossaryentry{avilissement}{
	name=avilissement,
	description={Action de se déprécier, de perdre de sa valeur; résultat de cette action.
	}
}
\newglossaryentry{double dépense}{
	name=double dépense,
	description={Le problème de double dépense renvoie, en informatique des systèmes décentralisés, à une situation ou un même débiteur serait en position de signer en envoyer sans intervalle de temps, deux transactions prenant pour point d'entrée une unique unité de monnaie-cryptographique, à destination de deux bénéficiaires distincts. En conséquence, la probabilité pour que deux ensembles de \emph{n} n\oe{}uds connaissent d'une information divergente sur l'état de la chaîne de transactions pendante accentue le risque de désaccord sur la validité du registre à l'échelle du réseau.
	}
}
\newglossaryentry{système pair à pair}{
	name=système pair à pair,
	description={Décrit un modèle d'échange en réseau où chaque entité est à la fois client et serveur. Dans le cadre du Bitcoin, ce système peut être qualifié de parfaitement décentralisé du fait que les connexions entre participants se réalisent sans intervention d'une infrastructure de répartition (tiers de confiance). Le partage d'information, du fait de la dynamique réticulaire, gagne en performance et en immutabilité : une censure ou tentative de modification des transaction ne sera efficiente qu'à condition de suspendre l'interconnexion entre les n\oe{}uds.
	}
}
\newglossaryentry{maximum supply}{
	name=maximum supply,
	description={Indicateur économique de la quantité maximale théorique d'émission d'un jeton de monnaie-cryptographique.
	}
}
\newglossaryentry{halving}{
	name=halving,
	description={Phénomène périodique de réduction de moitié de la prime de résultat associée au processus de validation d'un bloc de transactions du DEEP. Le halving s'interprète en un mécanisme de convergence de la masse monétaire vers l'offre maximale théorique imprimée lors du développement du protocole.
	}
}
\newglossaryentry{Livre blanc}{
	name=Livre blanc,
	description={Document d'information concis traitant des objectifs et moyens mis en \oe{}uvre à la production d'un dispositif d'enregistrement électronique partagé. Dans le cadre du règlement MiCA, ce dernier devant notamment contenir une description détaillée de l'émetteur, des droits, obligations et risques attachés aux crypto-actifs, ainsi qu'une présentation des principaux participants à la conception et à la mise au point du projet.
	}
}
